\documentclass[11pt]{article}
\usepackage[margin=2.5cm]{geometry}
\usepackage{amsmath,amssymb}
\usepackage{booktabs}
\usepackage{hyperref}
\usepackage{siunitx}

\newcommand{\Kone}{K_1}
\newcommand{\Ktwo}{K_2}
\newcommand{\etaone}{\eta_1}
\newcommand{\etatwo}{\eta_2}

\title{Identification of Twin-Related Grain Pairs from EBSD Orientation Data}
\author{}
\date{}

\begin{document}
\maketitle

\section{Problem statement}

Given two crystallographically indexed grains (clusters) with average
orientation matrices $G_1, G_2$ (mapping sample coordinates to crystal
coordinates), we wish to determine whether the pair is related by a
specific twinning mode, defined by its crystallographic elements
$\Kone$, $\Ktwo$, $\etaone$, $\etatwo$ (composition planes and shear
directions) tabulated in a twin-system database for a set of $M$
crystallographically equivalent variants of the mode.

The variant count $M$ is not a free parameter: it is the size of the
orbit of the twin-element tuple $(\Kone,\etaone)$ (or, more generally,
the full tuple $(\Kone,\Ktwo,\etaone,\etatwo)$, depending on how the
tabulation was generated) under the action of the phase's point group
$\mathcal G$, of order $N_s$. By the orbit--stabiliser theorem,
\begin{equation}
  M = \frac{N_s}{|\mathrm{Stab}(\Kone,\etaone)|},
\end{equation}
so $M \le N_s$ always, with equality only when the stabiliser of the
twin-element tuple is trivial. Different twin systems within the same
point group can therefore have different $M$; $M$ must be taken as
tabulated for each system rather than assumed. A concrete value of $M$
is given for the worked example in Section~\ref{sec:example}.

\section{Symmetry ambiguity of measured orientations}

A crystal orientation is only ever known modulo the crystal's point group
$\mathcal{G}$ (order $N_s$): for any $S \in \mathcal{G}$, $S G$ represents
the same physical crystal. Consequently, if $G_1$ and $G_2$ are each
independently resolved to \emph{some} representative (e.g.\ by an
unconstrained least-squares average over the grain's pixels), the two
representatives carry no guarantee of being labelled on
\emph{mutually consistent} branches. The raw misorientation
\begin{equation}
  M = G_2 G_1^{\mathsf T}
\end{equation}
is then only \emph{one arbitrary member} of the physically equivalent set
$\{\, S_2 G_2 (S_1 G_1)^{\mathsf T} : S_1,S_2 \in \mathcal G \,\}$, and its
rotation angle need not be close to the true minimum disorientation
between the two grains.

\paragraph{Resolution.} Before any twin analysis, every cluster's
average orientation (and, for consistency, its constituent pixel
orientations) is re-expressed on a single symmetry branch chosen so as to
minimise disorientation to one shared reference grain. This is a
\emph{single} symmetry operation per cluster, applied uniformly to that
cluster's average and to every one of its pixels -- never resolved
independently per pixel, since an independent per-pixel search was found
empirically to split a single physically coherent grain's pixels across
several different branches when pixel-level measurement noise straddles
a branch boundary. A concrete instance of both this failure mode and its
fix is given in Section~\ref{sec:example}.

\section{Twin conditions used}

A twinning operation can be described, for a Type~I twin, as a
$\SI{180}{\degree}$ rotation about the $\Kone$ plane normal $n_1$
(a reflection across $\Kone$ is experimentally indistinguishable from
this rotation for a centrosymmetric crystal, by Friedel's law -- EBSD
diffraction patterns cannot distinguish a crystal from its point-inverted
enantiomorph, so the effective symmetry group used throughout is the
extended (Laue-class) group, improper operations included). For a
Type~II twin, the equivalent description is a $\SI{180}{\degree}$
rotation about the $\etaone$ shear direction $a_1$. A third, general
description (valid for compound and non-compound modes alike) is a
rotation about the $\Ktwo$ plane normal $n_2$ by the mode's
characteristic shear angle $\gamma$.

All three conditions are checked:
\begin{align}
  \text{K1\_180:}\quad & \text{axis}(M) \parallel n_1, \quad \text{angle}(M) = \SI{180}{\degree} \\
  \text{eta1\_180:}\quad & \text{axis}(M) \parallel a_1, \quad \text{angle}(M) = \SI{180}{\degree} \\
  \text{K2\_shear:}\quad & \text{axis}(M) \parallel n_2, \quad \text{angle}(M) = \gamma
\end{align}
where $\text{axis}(M)$, $\text{angle}(M)$ are extracted from $M$ via a
sign-consistent unit quaternion, and axis comparison is sign-independent
($|\hat n \cdot \hat a|$), since a rotation axis and the tabulated
crystallographic directions are each defined only up to line direction.

All three conditions are always fully evaluated for every candidate pair
(never short-circuited on the first match), and every condition
satisfying the tolerance is reported, not only the single best. A pair
for which K1\_180 and eta1\_180 are \emph{both} independently satisfied
(generally via different representative rotations, see below) is a
compound twin -- all three descriptions coincide exactly when the mode's
$\Kone,\Ktwo,\etaone,\etatwo$ are all simultaneously rational; observing
K1\_180, eta1\_180, and K2\_shear simultaneously satisfied for the same
pair is accordingly the expected, confirmatory signature of a compound
twin.

\section{The representative-rotation problem}

A subtlety was discovered during validation that is easy to miss: even
after the branch-consistency fix above, comparing the \emph{minimum-angle}
representative of $M$ against the $\SI{180}{\degree}$ conditions can fail
to find a genuine Type~I/II relationship. This is because conjugation,
$M \mapsto S M S^{\mathsf T}$, preserves rotation angle exactly and only
rotates the axis; if the true minimum disorientation angle is far from
$\SI{180}{\degree}$ (as it generically is), no amount of searching over
conjugation alone will ever produce a $\SI{180}{\degree}$ candidate.

The complete search requires a second, non-conjugation degree of freedom,
distinct from the target-side variant search introduced above: for every
proper $T \in \mathcal G$,
\begin{equation}
  M_T = M\,T
\end{equation}
is tested against the tabulated target directions (each already
representing $M_{\text{var}}$ equivalent variants -- using $M_{\text{var}}$
here to avoid clashing with the misorientation symbol $M$ -- so that
testing $M_T$ against all $M_{\text{var}}$ stored variants is equivalent
to an additional conjugation search on the target side). Algebraically,
$M_T = G_2 (T^{\mathsf T} G_1)^{\mathsf T}$: this is exactly the same
construction as $M$, but with $G_1$ replaced by the equally-valid
symmetric branch $T^{\mathsf T} G_1$. The $T$-loop combined with the
target-variant search was confirmed, by direct comparison against a
brute-force search over the complete two-sided set $\{S_2 M
S_1^{\mathsf T} : S_1,S_2\in\mathcal G\}$, to reproduce the true global
optimum exactly. A worked instance is given in Section~\ref{sec:example}.

\section{Algorithm summary}

For a candidate grain pair with (branch-consistent) average orientations
$G_1, G_2$:
\begin{enumerate}
  \item Form $M = G_2 G_1^{\mathsf T}$.
  \item For every proper symmetry operation $T$ of the phase:
  \begin{enumerate}
    \item Form $M_T = M T$ and extract its rotation axis and angle.
    \item For each of the three target families ($n_1$, $a_1$, $n_2$;
      each with $M_{\text{var}}$ tabulated equivalent variants): find
      the variant minimising the combined axis and angle deviation from
      the condition's target axis and target angle
      ($\SI{180}{\degree}$ for K1\_180/eta1\_180, $\gamma$ for K2\_shear).
  \end{enumerate}
  \item For each condition, retain the globally best $(T, \text{variant})$
    pair found over the full search.
  \item Report the pair as a match for a given condition if both the
    axis and angle deviations fall below a tolerance
    $\text{tol}=\SI{5}{\degree}$ (default); report \emph{all} conditions
    satisfying this, not only the single best.
\end{enumerate}

For $N$ candidate grains this is applied to every unordered pair
($\binom{N}{2}$ total), vectorised across pairs, symmetry operations, and
target variants. The cost per pair scales as $O(N_s \times
M_{\text{var}})$ per condition.

\paragraph{Relation to a naive brute-force search.} A brute-force search
of the complete two-sided symmetric-equivalence set against a single
fixed target vector, without using the tabulated variants at all, would
require looping $S_1$ and $S_2$ independently over the full point group,
i.e.\ $O(N_s^2)$ candidate evaluations. Looping only over the $M$
distinct tabulated variants instead of over all $N_s$ raw symmetry
operations on the target side is exact and lossless, not an
approximation: by construction the $M$ variants are precisely the
distinct images produced by the $N_s$ operations acting on the twin
element tuple, so iterating the $M$ pre-computed variants reproduces the
identical set of comparisons that iterating all $N_s$ raw operations
would produce, with the redundant (duplicate-image) evaluations already
removed. The $T$-loop-plus-variant-search cost is therefore smaller than
the naive $O(N_s^2)$ brute force by a factor
\begin{equation}
  \frac{N_s^2}{N_s \, M_{\text{var}}} = \frac{N_s}{M_{\text{var}}}
  = |\mathrm{Stab}(\Kone,\etaone)|,
\end{equation}
i.e.\ exactly the stabiliser order of the twin-element tuple. Twin
systems with a small number of unique variants $M_{\text{var}}$ (a large
stabiliser) are therefore searched proportionally faster; the search is
never more expensive than, and is equal to, the brute-force cost only in
the worst case $M_{\text{var}} = N_s$ (trivial stabiliser).

\section{Worked example: the \{114\} austenite twinning mode in NiTi}
\label{sec:example}

The general method above is illustrated with the $\{114\}$ deformation
twin mode in B2 (austenite) NiTi, generated from the tabulated primary
variant
\begin{equation}
  \etaone = [1\,2\,2],\qquad \Kone = (\bar 4\,1\,1),
\end{equation}
expressed in the crystal's real-space and reciprocal-space bases
respectively, converted to Cartesian unit vectors $a_1, n_1$ via the
phase's structure matrices. The remaining $M-1$ equivalent variants are
generated by applying every proper cubic symmetry operation
($N_s = 24$) to the tuple $(\Kone,\Ktwo,\etaone,\etatwo)$ and retaining
only genuinely distinct results -- two candidates are treated as
equivalent if \emph{either} their $\Kone$ \emph{or} their $\etaone$
coincide (up to sign) with an already-retained variant. This is a
coarser equivalence than a single-vector orbit (which, checked directly
for $n_1$ alone, has orbit size $3$ under the $24$ proper rotations, not
$M$), consistent with the general bound $M \le N_s$ derived above; for
this system the generation algorithm yields $M = 12$, i.e.\ a search
cost reduction of a factor $N_s/M_{\text{var}} = 24/12 = 2$ relative to
the naive $O(N_s^2)$ brute-force search, per the general result above.

\paragraph{Branch-consistency failure, observed.} On real EBSD data, one
grain pair's independently-resolved (unsymmetrized) average orientations
gave a raw misorientation angle of \SI{68.66}{\degree}; after resolving
both grains to branches consistent with a shared reference, the same
pair's misorientation was \SI{37.00}{\degree}, matching the true minimum
disorientation and the value visible directly on the pair's measured
pole figures. Separately, for one cluster, independent per-pixel branch
resolution (rather than the single-symmetry-operation-per-cluster
approach specified above) was found to split its $1408$ pixels across
$6$ distinct branches, despite an internal pixel-to-pixel orientation
spread of only $\approx\SI{6}{\degree}$ once resolved consistently --
confirming the branch ambiguity is a real, non-negligible effect on this
class of data, not a theoretical concern only.

\paragraph{Representative-rotation search, observed.} For a confirmed
grain pair, the minimum-angle representative of $M$ had rotation angle
$\approx\SI{37}{\degree}$, matching the K2\_shear description (axis
deviation \SI{1.28}{\degree}, angle deviation \SI{1.83}{\degree}) -- but
this angle is nowhere near \SI{180}{\degree}, so neither K1\_180 nor
eta1\_180 could match via conjugation alone at this representative. The
$T$-loop search found a specific alternative representative matching
K1\_180 with axis deviation \SI{1.00}{\degree} and angle deviation
\SI{0.15}{\degree} -- an essentially exact Type~I twin relationship,
confirmed independently by brute-force search over the complete
two-sided symmetric-equivalence set, and consistent with the pair's
directly observed coincident $\{114\}$ pole on the measured pole
figures. A second, independently checked grain pair was found to satisfy
K1\_180, eta1\_180, \emph{and} K2\_shear simultaneously, each via its own
representative rotation $T$ -- the expected signature of a compound
twin, per Section~3.

\section{Tolerance and interpretation}

Because the search evaluates, per pair and per condition, $N_s \times
M_{\text{var}}$ candidate rotations (of order several hundred), a
nominal tolerance tuned for a much smaller search (e.g.\ the original
target-variant-only search) is not automatically appropriate once the
$T$-loop is included: the enlarged search finds the true global optimum
rather than an artifact of incomplete search, but the achievable minimum
deviation under \emph{no} true relationship also shrinks as the number
of candidates grows. The tolerance should therefore be calibrated
against the specific dataset's orientation-measurement noise floor
(typically a few degrees for EBSD) rather than treated as a universal
constant, and the reported deviation values should be inspected directly
(e.g.\ sorted by deviation) rather than relying solely on a hard
pass/fail cutoff.

\end{document}
